% chapters/15/evaluasi.tex
\bagian{Asesmen dan Instrumen Penilaian}
\label{sec:asesmen-audio}

\begin{asesmenbox}
\textbf{Kedudukan asesmen bab ini.} Bab ini menuntaskan \SUBCPMK{4.2}, yaitu
bab kedua pada \CPMK{4}. Ketiga indikatornya diperiksa melalui penilaian
artefak \istilah{project} (Tugas Praktikum 12) dan soal konsep.

\vspace{0.4em}
Materi Bab~XV termasuk bahan \textbf{Proyek Akhir} pada minggu kelima belas
dan \textbf{Ujian Akhir Semester} pada minggu keenam belas.
\end{asesmenbox}

\subsection*{Instrumen 1: Rubrik Praktikum 12}

\begin{table}[htbp]
\centering
\caption{Rubrik penilaian Praktikum 12 beserta indikator yang diukur}
\label{tab:rubrik-praktikum-12}
{\setlength{\tabcolsep}{4pt}\footnotesize%
\begin{tabularx}{\textwidth}{@{}L{3.2cm} C{1.3cm} C{2.1cm} Y@{}}
\toprule
\headrow \textbf{Kriteria} & \textbf{Bobot} & \textbf{Indikator} & \textbf{Deskripsi capaian} \\
\midrule
Peristiwa dan pemicu bunyi & 40\% & \ind{a} &
Setiap peristiwa pada daftar memicu bunyi yang benar, pada saat yang tepat, dan
bunyi yang berulang tidak berbunyi setiap bingkai. \\
\addlinespace
Pemilihan jenis bunyi & 25\% & \ind{b} &
Musik latar dan efek suara dipisahkan dengan benar, baik pada sumber bunyinya
maupun pada tingkat \istilah{AudioMixer}, dan alasannya dijelaskan. \\
\addlinespace
Waktu dan kekuatan bunyi & 25\% & \ind{c} &
Tingkat kekuatan tersusun dan beralasan; efek suara terdengar jelas di atas
musik latar; hasil keempat percobaan tercatat. \\
\addlinespace
Partikel dan catatan aset & 10\% & \ind{a}, \ind{c} &
Sekurang-kurangnya satu partikel dipasang pada peristiwa yang tepat, dan daftar
berkas bunyi memuat sumber serta lisensinya. \\
\bottomrule
\end{tabularx}}
\end{table}

\subsection*{Instrumen 2: Soal Konsep}

Jawablah dengan singkat. Kunci jawaban beserta indikator yang diukur disajikan
pada Tabel~\ref{tab:kunci-konsep-ch15}.

\begin{enumerate}
  \item Jelaskan perbedaan \code{AudioClip} dan \code{AudioSource}.
  \item Sebuah \code{AudioSource} dipakai untuk musik dan efek suara sekaligus.
        Gejala apa yang muncul?
  \item Berapa besar perubahan kekuatan bunyi apabila \code{sfxVolume} dinaikkan
        dari 0,25 menjadi 1,0? Nyatakan dalam desibel.
  \item Musik bernilai $-6$ dB dan efek suara bernilai $-6$ dB. Mengapa bunyi
        lompatan tetap terdengar samar, dan apa yang sebaiknya diubah?
  \item Bunyi langkah dipanggil berulang dengan cepat dan terdengar seperti
        mesin. Sebutkan dua cara memperbaikinya.
\end{enumerate}

\begin{table}[htbp]
\centering
\caption{Kunci soal konsep Bab XV beserta indikator yang diukur}
\label{tab:kunci-konsep-ch15}
{\setlength{\tabcolsep}{4pt}\footnotesize%
\begin{tabularx}{\textwidth}{@{}C{0.8cm} C{2.1cm} Y@{}}
\toprule
\headrow \textbf{No} & \textbf{Indikator} & \textbf{Pokok jawaban} \\
\midrule
1 & \ind{b} &
\code{AudioClip} adalah berkas bunyinya, yaitu data yang tidak dapat berbunyi
sendiri dan tidak memiliki letak. \code{AudioSource} adalah komponen pada
\istilah{GameObject} yang memutar bunyi itu dan menentukan letaknya di dalam
ruang. \\
\addlinespace
2 & \ind{b} &
Setiap efek suara menghentikan musik yang sedang berjalan, sebab
\code{Play} memakai satu sumber untuk satu bunyi pada satu waktu. Gejalanya:
musik terputus-putus setiap kali pemain melompat atau mengambil benda. \\
\addlinespace
3 & \ind{a} &
Dari 0,25 ke 1,0 berarti kekuatannya menjadi empat kali. Karena
$20 \times \log_{10} 4 \approx 12$, perubahannya sekitar 12 dB. \\
\addlinespace
4 & \ind{c} &
Sebab keduanya sama kuat, dan pada kekuatan yang sama telinga mendahulukan
bunyi yang berjalan terus-menerus. Yang diubah adalah musiknya: diturunkan
sekitar 10 dB agar bunyi lompatan memperoleh keunggulan yang cukup. Menaikkan
efek suaranya justru berisiko membuat bunyinya pecah. \\
\addlinespace
5 & \ind{c} &
Pertama, nada digeser sedikit dan berbeda-beda setiap kali dipanggil. Kedua,
beberapa berkas bunyi langkah dipakai bergantian, sehingga pengulangannya tidak
terdengar. \\
\bottomrule
\end{tabularx}}
\end{table}
